\documentclass[10pt,a4paper]{amsart}
\usepackage[margin=19mm]{geometry}
\usepackage{amsmath,amssymb,mathtools}
\usepackage[scr=rsfs,cal=euler]{mathalfa}
\usepackage{xfrac}
\usepackage[unicode,hidelinks,bookmarks=false]{hyperref}
\newcommand{\ie}{i.e.\ }
\newcommand{\cf}{cf.\ }
\newcommand{\resp}{respectively\ }
\newcommand{\Subsection}{\subsection}
\providecommand{\nobreakdash}{\nobreak\hbox{-}}
\setlength{\emergencystretch}{2em}
\multlinegap0pt
\usepackage{amssymb}
\usepackage{float}
\newtheorem{corollary}{Corollary}
\title{Towards a Theory of Pragmatic Information}
\author{Edward D. Weinberger}
\date{Source: arXiv:2403.12324v8 (CC BY 4.0)}
\begin{document}
\maketitle

\noindent\emph{Source and licence.} Towards a Theory of Pragmatic Information. Version arXiv:2403.12324v8 (CC BY 4.0), distributed by arXiv under the Creative Commons Attribution 4.0 International licence, http://creativecommons.org/licenses/by/4.0/, as recorded in the arXiv licence metadata for this exact version and shown on the abstract page https://arxiv.org/abs/2403.12324v8. Full licence notice: \texttt{licenses/arxiv-2403.12324-CC-BY-4.0.txt}.\par\smallskip

\noindent\textit{Editorial scope.} Counted unit: the article's corollary eq:PragIndependenceCond (source0.tex lines 458--470). The article's complete original proof of it is delivered with it (source0.tex lines 472--486, one \texttt{proof} environment of 15 lines). No mathematics is written, repaired or re-derived: the statement and the proof are the article's own lines, in the article's own notation, and no retained line is substituted or re-flowed.  No further source-local context is needed: every cross-reference of every retained line resolves inside the two delivered spans.  Nothing was omitted from any retained span, so the package carries no omission record.  No retained line contains a citation, so the wrapper carries no bibliography and no bibliography entry.  No retained line contains a figure environment, an image-inclusion command or drawing code, so no image is delivered and the package declares no asset dependency.  Editorial formatting only, in addition: an AMS \texttt{amsart} preamble that restates the article's own declarations and macros used by the retained lines, namely the environments \texttt{corollary} on separate counters, together with the standard packages those lines need.  The retained environments print in this document as Corollary 1 in order of appearance, whereas the article prints the same items with the numbering of its own full document.  The complete native source of the article as distributed in the arXiv e-print for this exact version is bundled unchanged as \texttt{sources/156301/source0.tex}.
\begin{corollary}[Additivity of Pragmatic Information]
If, in addition, to the hypotheses of  Theorem 3.2.4, we have
\begin{align}
q'_{i'|i} = q'_{i'} {\it \quad and \quad} p_{i, i', m, m'} = p_{i, m} p'_{i', m'}  {\it \ for \ all \ } i, i', m, m', 
\label{eq:PragIndependenceCond}
\end{align}
then 
\begin{align}
 \Phi_{\Delta, \Delta'} (\mathcal M, \mathcal M'; \Omega, \Omega')=\Phi_{\Delta} ({\mathcal M}; \Omega)  +
\Phi_{\Delta'} ({\mathcal M'}; \Omega').   
\label{eq:PragIndependence}
\end{align}
\end{corollary}

\begin{proof}
Summing over $i$ and $i'$, we conclude that $\varphi_{m, m'} = \varphi_m \varphi_{m'}$.  Thus, 
\begin{align*}
p_{i, i' | m, m'} &= \frac{p_{i, i', m, m'}}{\varphi_{m, m'}} \\
                       &= \left[\frac{p_{i, m}}{\varphi_{m}}\right] \left[\frac{p'_{i', m'}}{\varphi_{m'}}\right] \\
                       &= p_{i|m} p'_{i'|m'}.  \\
\end{align*}
The result follows, since
\begin{align*}
\Phi_{\Delta, \Delta'} (\mathcal M, \mathcal M'; \Omega, \Omega') 
                       &= \sum_{i, i', m, m'} p_{i, i', m, m'} \log \left(\frac{p_{i,i'|m, m'}}{q_{i,i'}}\right)  \\
                       &= \sum_{i, i', m, m'}p_{i, m} p'_{i', m'} \log \left(\frac{p_{i|m} p_{i'|m'}}{q_{i}q_{i'}}\right)  \\
                       &= \Phi_{\Delta} ({\mathcal M}; \Omega)  + \Phi_{\Delta'} ({\mathcal M'}; \Omega').\\
\end{align*}
\end{proof}

\end{document}
